\section{统一求解流程与策略落地}

\subsection{从预测到真实账单的双层流程}

四问共享“决策模型”和“执行评价”两层。决策层只读当时可用的信息，输出正常购电和电池计划；评价层才读入当天真实负荷、光伏与价格，逐10 min回放并计算紧急购电。若将评价层的真实量送回同日更早的决策，就会产生不可执行的信息泄漏。本文以时间戳断言和扰动测试同时检查这一边界。

\begin{figure}[H]
\centering
\begin{tikzpicture}[node distance=8mm and 10mm,
box/.style={draw,rounded corners,align=center,minimum height=9mm,text width=3.0cm},
arr/.style={-{Latex[length=2mm]},thick}]
\node[box] (hist) {历史完整日\\已发布预报};
\node[box,right=of hist] (fc) {点预测与\\条件场景};
\node[box,right=of fc] (opt) {随机优化\\输出计划};
\node[box,below=of opt] (exec) {仅执行到\\下一更新点};
\node[box,left=of exec] (real) {真实观测到达\\更新SOC};
\node[box,left=of real] (bill) {真实样本回放\\独立重算账单};
\draw[arr] (hist)--(fc);
\draw[arr] (fc)--(opt);
\draw[arr] (opt)--(exec);
\draw[arr] (exec)--(real);
\draw[arr] (real)--(bill);
\draw[arr,dashed] (real)--(fc);
\node[draw=red!70,dashed,fit=(hist)(fc)(opt),inner sep=4mm,label=above:{\small 决策时可用信息}] {};
\node[draw=blue!60,dashed,fit=(exec)(real)(bill),inner sep=4mm,label=below:{\small 执行与样本外评价}] {};
\end{tikzpicture}
\caption{统一的因果决策、滚动执行和真实账单复核流程}
\label{fig:algorithm-flow}
\end{figure}

对问题二，每个目标日执行以下步骤：首先用目标日前历史更新候选预测器；其次从最近28个样本外残差日抽取并压缩场景；再次求解共享日前变量的两阶段模型；最后以当天真实值计算补救量，保存日末库存并传给下一日。问题三与问题四的日内调整版本在一天内重复“更新信息集—重建剩余场景—求解—执行一个区间”，已经执行的决策冻结不变。该滚动实现比一次性建立全年多阶段场景树更易审计，也避免场景数随更新次数指数增长。

\subsection{参数选择与防止事后挑选}

所有影响主策略的离散选择先在1月验证期冻结，2月至12月仅做一次连续样本外评价。表\ref{tab:fixed-parameters}区分物理参数、题面结算参数和数据驱动参数；前两类不参与调优，后一类才由历史验证选择。正式期中表现更好的备选仍保留在消融表中，但不会替换主方案。

\begin{table}[H]
\centering
\caption{关键参数来源、固定值与作用}
\label{tab:fixed-parameters}
\begin{tabularx}{\textwidth}{p{2.5cm}p{3.0cm}p{2.5cm}X}
\toprule
类别 & 参数 & 主值 & 依据及作用\\
\midrule
物理 & $\Delta t,E_{\min},E_{\max}$ & $1/6$ h，1200，10800 kWh & 题面给定，统一控制时间与库存边界\\
物理 & $\eta_c,\eta_d,\bar P$ & 0.9，0.9，5000 kW & 主假设与题面功率上限\\
结算 & 紧急购电倍率 & 5 & 题面给定，进入真实现金账单\\
结算 & 调整系数 & 上调1.5、下调0.5 & 以原计划为基准一次结算\\
问题二 & 预测器、场景数 & 随机森林、256抽取 & 1月滚动费用验证；重复场景压缩赋权\\
问题二 & 终端倍率$m$ & 1.2 & 近似次日库存价值\\
问题三 & 历史窗、场景数 & 28日、64抽取 & 校准提前量误差并控制计算规模\\
问题三 & 风险参数 & $\alpha=0.9,\lambda=0.1$ & 作为风险增强备选，不替代风险中性主方案\\
问题四 & 价格基线 & 28日均值 & 1月历史验证优于备选基线\\
问题四 & 修正与终端参数 & $\alpha_p=0.2$，衰减0.1，$\gamma=0.5$ & 冻结日内修正速度与日末库存价值\\
问题四 & $\rho,K$ & $0,1$ & 1月验证选择；$K=2,3$仅作敏感性\\
\bottomrule
\end{tabularx}
\end{table}

这种“验证一次、正式期锁定”的设计牺牲了事后最优数字，却使方案在真实运行时可复现。尤其是第四问$K=3$的正式期费用较低，但若根据这一结果改选，就等价于让未来334天参与模型选择；本文明确拒绝这一做法。

\subsection{算法规模、输出和异常处理}

求解器输入全部使用kWh和元，外部展示需要功率时再除以$\Delta t$。每次求解后先检查状态方程、母线平衡、变量上下界和互斥或净化条件，再把执行区间写入轨迹。日末执行以下守卫：库存不得因四舍五入漂移，下一日直接读取原浮点值；最后一天施加6000 kWh终端条件；缺少历史滞后或场景无法分组时回退到预先规定的简单模型，而非读取未来补齐。

\begin{table}[H]
\centering
\caption{正式计算规模与独立复核证据}
\label{tab:computation-audit}
\begin{tabularx}{\textwidth}{p{2.1cm}rrX}
\toprule
任务 & 正式日数 & 求解或记录规模 & 主要复核结论\\
\midrule
问题一 & 1 & LP与MILP各1次 & 目标一致，MILP严格互斥\\
问题二 & 334 & 25组实验、4837条日模型/回放记录 & 28项测试通过，最大残差$3.73\times10^{-9}$\\
问题三 & 334 & 21组实验、28056次滚动求解 & 20项测试通过，最大残差$7.55\times10^{-7}$\\
问题四 & 334 & 29组、3912策略日、10248次求解 & 18项测试通过，最大误差$1.863\times10^{-9}$\\
\bottomrule
\end{tabularx}
\end{table}

最终输出分为三层：论文只保留关键模型、主结果和足以解释结论的图表；结果工作簿保存题目要求的指定时段、日期汇总与完整轨迹；附录和支撑材料保存全部源程序、配置、测试与审计记录。图表由正式CSV或工作簿自动生成，未手工改写数据点。这样既保持正文可读，也能从任一汇总数追溯到10 min执行记录。

\subsection{跨问题结果的可比口径}

表\ref{tab:cross-question}汇总各问主结果，但不同问题的价格、信息和结算权限不同，绝不能横向把总费差额直接解释为某个算法的收益。有效比较只发生在同一问题内部、使用相同正式日期与评价账单的对照之间。

\begin{table}[H]
\centering
\caption{四问主方案结果及正确比较口径}
\label{tab:cross-question}
\begin{tabularx}{\textwidth}{p{1.3cm}p{3.2cm}rrX}
\toprule
问题 & 主方案 & 总费/万元 & 紧急量/MWh & 可支持的结论\\
\midrule
一 & 确定性MILP & 3.51 & 0 & 相对无储能同日基线下降26.898\%\\
二 & 随机森林+联合场景 & 1447.93 & 262.35 & 相对同预测器点方案下降7.690\%\\
三 & 条件尺度随机MPC & 1662.77 & 238.47 & 四次新预报相对仅0:00新预报节省41.21万元\\
四（日前） & 边际独立收缩 & 1520.11 & 256.56 & 相对波动价点方案下降7.576\%\\
四（日内） & 一组条件前瞻 & 1499.72 & 176.24 & 相对同机制点方案下降6.605\%\\
\bottomrule
\end{tabularx}
\end{table}

问题二与问题三使用题面固定电价，问题四使用真实波动价；问题三和问题四日内版本允许调整正常购电，而问题二不允许。因此，跨行比较只能用于描述任务递进，不能据此判断哪个模型“绝对更优”。这一限制与样本外显著性区间一起构成本文结论边界。
